\centerline{\bf \Huge Домашнее задание.}
\centerline{\bf \Huge Рыцари и лжецы. Уравнения.}

\begin{flushright}
        \scriptsize{Пусть я никогда не видел его победителем,\\ но и сдавшимся после поражения тоже. \\Ханнес. Атака Титанов}
\end{flushright}

\problem{1}В ЭлМШ обучается 41 ученик: шиншиллы, лососи и киви. Шиншиллы на все вопросы отвечают правильно, киви всегда ошибаются, а лососи на заданные им вопросы строго по очереди отвечают то верно, то ошибаются. Всем ученикам ЭлМШ было задано по три вопроса: «Ты шиншилла?», «Ты лосось?», «Ты киви?». Ответили «Да» на первый вопрос — 25 учащихся, на второй — 14, на третий — 2. Сколько лососей учатся в ЭлМШ?

\problem{2}В лицее каждый восьмиклассник либо физмат, который всегда говорит правду, либо гум, который всегда лжет. Елена Михайловна спросила каждого восьмиклассника про каждого из остальных восьмиклассников, физмат тот или гум. Всего было получено 144 ответа «физмат» и 96 ответов «гум». Сколько правдивых ответов могла получить Елена Михайловна?

\problem{$3^{*}$} Физматы и гумы решил зарыть топор войны и сели за круглый стол переговоров. Известно, что суммарно на переговоры пришло 12 учеников, физматы всегда говорят правду, гумы всегда лгут. Елена Михайловна рассадила ребят за круглым столом так, что они образуют правильный двенадцатиугольник. Ребята забыли представиться, кто из них физмат, а кто гум, поэтому Елене Михайловне пришлось разбираться самой. Она может встать куда угодно и спросить у всех учеников: "Каково расстояние от меня до ближайшего гуманитария?". После этого каждый ей отвечает на этот вопрос. Какое минимальное количество вопросов должна задать Елена Михайловна, чтобы гарантировано узнать, кто из них физмат, а кто гум. (Мероприятие переговоров приличное и конфдинциальное, поэтому на стол вставать нельзя и посторонних людей, кроме учеников и Елены Михайловны там нет.) \\
\textit{P.s. Привет от Расула Владимировича :)}